\begin{ex}
Tính giá trị của biểu thức: $4 + 5 - 3$.
\shortans{6}
\loigiai{
Thực hiện từ trái sang phải: $4 + 5 = 9$; sau đó $9 - 3 = 6$.
\textbf{Đáp số:} $6$.
}
\end{ex}

\begin{ex}
Tính giá trị của biểu thức: $10 - 6 + 4$.
\shortans{8}
\loigiai{
Thực hiện từ trái sang phải: $10 - 6 = 4$; sau đó $4 + 4 = 8$.
\textbf{Đáp số:} $8$.
}
\end{ex}

\begin{ex}
Tính giá trị của biểu thức: $8 - 5 + 7$.
\shortans{10}
\loigiai{
Thực hiện từ trái sang phải: $8 - 5 = 3$; sau đó $3 + 7 = 10$.
\textbf{Đáp số:} $10$.
}
\end{ex}

\begin{ex}
Điền số thích hợp vào ô trống: $3 + \dots - 2 = 6$.
\shortans{5}
\loigiai{
Ta có: $3 + \dots = 6 + 2 = 8 \Rightarrow$ số cần điền là: $8 - 3 = 5$.
\textbf{Đáp số:} $5$.
}
\end{ex}

\begin{ex}
Điền số thích hợp vào ô trống: $10 - \dots + 3 = 8$.
\shortans{5}
\loigiai{
Ta có: $10 - \dots = 8 - 3 = 5 \Rightarrow$ số cần điền là: $10 - 5 = 5$.
\textbf{Đáp số:} $5$.
}
\end{ex}

\begin{ex}
Tìm số tự nhiên $x$ thỏa mãn: $x + 4 - 1 = 8$.
\shortans{5}
\loigiai{
Ta có: $x + 3 = 8 \Rightarrow x = 8 - 3 = 5$.
\textbf{Đáp số:} $5$.
}
\end{ex}

\begin{ex}
Tìm số tự nhiên $y$ thỏa mãn: $9 - y + 2 = 7$.
\shortans{4}
\loigiai{
Ta có: $9 - y = 7 - 2 = 5 \Rightarrow y = 9 - 5 = 4$.
\textbf{Đáp số:} $4$.
}
\end{ex}

\begin{ex}
Lan có 5 cái nhãn vở. Mẹ cho Lan thêm 4 cái, Lan dùng hết 3 cái. Hỏi Lan còn lại bao nhiêu cái nhãn vở?
\shortans{6}
\loigiai{
Số nhãn vở Lan có sau khi mẹ cho là: $5 + 4 = 9$ (cái).\nSau khi dùng 3 cái, Lan còn lại: $9 - 3 = 6$ (cái nhãn vở).
\textbf{Đáp số:} $6$.
}
\end{ex}

\begin{ex}
Bác thợ mộc có một thanh gỗ dài 10 cm. Bác cưa bớt đi 6 cm, sau đó đóng nối thêm một đoạn gỗ 5 cm. Hỏi thanh gỗ mới dài bao nhiêu xăng-ti-mét?
\shortans{9}
\loigiai{
Chiều dài thanh gỗ mới là: $10 - 6 + 5 = 9$ (cm).
\textbf{Đáp số:} $9$.
}
\end{ex}

\begin{ex}
Đàn chim có 7 con đậu trên cây. Có 3 con bay đi, lát sau có 5 con chim khác bay đến đậu. Hỏi trên cây lúc này có bao nhiêu con chim?
\shortans{9}
\loigiai{
Số chim còn lại sau khi 3 con bay đi là: $7 - 3 = 4$ (con).\nSố chim trên cây lúc này là: $4 + 5 = 9$ (con chim).
\textbf{Đáp số:} $9$.
}
\end{ex}

\begin{ex}
Điền dấu thích hợp ($>, <, =$): $4 + 6 \dots 3 + 7$. Kết quả của mỗi vế bằng bao nhiêu?
\shortans{10}
\loigiai{
Vế trái: $4 + 6 = 10$. Vế phải: $3 + 7 = 10$. Cả hai vế đều bằng $10$.
\textbf{Đáp số:} $10$.
}
\end{ex}

\begin{ex}
So sánh hai biểu thức: $10 - 4$ và $2 + 5$. Giá trị của biểu thức lớn hơn là bao nhiêu?
\shortans{7}
\loigiai{
Biểu thức 1: $10 - 4 = 6$.\nBiểu thức 2: $2 + 5 = 7$.\nVì $7 > 6$ nên giá trị của biểu thức lớn hơn là $7$.
\textbf{Đáp số:} $7$.
}
\end{ex}

\begin{ex}
Tìm số tự nhiên lớn nhất có một chữ số thỏa mãn: số đó trừ đi 3 có kết quả bé hơn 5.
\shortans{7}
\loigiai{
Ta có: $7 - 3 = 4 < 5$. Nếu là số 8: $8 - 3 = 5$ (không bé hơn 5).\nVậy số lớn nhất thỏa mãn là số $7$.
\textbf{Đáp số:} $7$.
}
\end{ex}

\begin{ex}
Tìm một số biết rằng khi lấy số đó cộng với 3 rồi trừ đi 4 thì được 5.
\shortans{6}
\loigiai{
Số cần tìm là: $5 + 4 - 3 = 6$.
\textbf{Đáp số:} $6$.
}
\end{ex}

\begin{ex}
Tìm một số biết rằng khi lấy 10 trừ đi số đó rồi cộng thêm 2 thì được 9.
\shortans{3}
\loigiai{
Ta có: $10 - x = 9 - 2 = 7 \Rightarrow x = 10 - 7 = 3$.
\textbf{Đáp số:} $3$.
}
\end{ex}

\begin{ex}
Trong một giỏ đựng 6 quả táo và 4 quả lê. Nam ăn 2 quả táo và 1 quả lê. Hỏi trong giỏ còn lại bao nhiêu quả cả hai loại?
\shortans{7}
\loigiai{
Tổng số quả ban đầu là: $6 + 4 = 10$ (quả).\nTổng số quả Nam đã ăn là: $2 + 1 = 3$ (quả).\nSố quả còn lại là: $10 - 3 = 7$ (quả).
\textbf{Đáp số:} $7$.
}
\end{ex}

\begin{ex}
Có 10 que tính. Bình dùng 3 que xếp 1 hình tam giác, dùng 4 que xếp 1 hình vuông. Hỏi Bình còn thừa bao nhiêu que tính?
\shortans{3}
\loigiai{
Tổng số que tính Bình đã dùng là: $3 + 4 = 7$ (que tính).\nSố que tính còn thừa là: $10 - 7 = 3$ (que tính).
\textbf{Đáp số:} $3$.
}
\end{ex}

\begin{ex}
Mai có 8 viên kẹo, Mai cho An 2 viên, sau đó mẹ lại cho Mai 3 viên nữa. Hỏi Mai có tất cả bao nhiêu viên kẹo?
\shortans{9}
\loigiai{
Số kẹo Mai còn lại sau khi cho An là: $8 - 2 = 6$ (viên kẹo).\nSố kẹo Mai có sau khi mẹ cho thêm là: $6 + 3 = 9$ (viên kẹo).
\textbf{Đáp số:} $9$.
}
\end{ex}

\begin{ex}
Tính: $2 + 3 + 5 - 4$.
\shortans{6}
\loigiai{
Thực hiện lần lượt: $2 + 3 = 5$; $5 + 5 = 10$; $10 - 4 = 6$.
\textbf{Đáp số:} $6$.
}
\end{ex}

\begin{ex}
Tính: $10 - 2 - 3 + 1$.
\shortans{6}
\loigiai{
Thực hiện lần lượt: $10 - 2 = 8$; $8 - 3 = 5$; $5 + 1 = 6$.
\textbf{Đáp số:} $6$.
}
\end{ex}

\begin{ex}
Lớp 1C có 10 bạn tham gia trò chơi. Có 4 bạn rút khỏi cuộc chơi, nhưng lại có thêm 2 bạn khác xin vào chơi. Hỏi lúc này có bao nhiêu bạn đang chơi?
\shortans{8}
\loigiai{
Số bạn đang tham gia chơi là: $10 - 4 + 2 = 8$ (bạn).
\textbf{Đáp số:} $8$.
}
\end{ex}

\begin{ex}
Tìm số tự nhiên $n$ biết rằng: $5 + n = 10 - 2$.
\shortans{3}
\loigiai{
Vế phải: $10 - 2 = 8$.\nTa có: $5 + n = 8 \Rightarrow n = 8 - 5 = 3$.
\textbf{Đáp số:} $3$.
}
\end{ex}

\begin{ex}
Tìm số tự nhiên $k$ biết rằng: $k - 1 = 4 + 4$.
\shortans{9}
\loigiai{
Vế phải: $4 + 4 = 8$.\nTa có: $k - 1 = 8 \Rightarrow k = 8 + 1 = 9$.
\textbf{Đáp số:} $9$.
}
\end{ex}

\begin{ex}
Một bể nước có 8 lít nước. Người ta múc ra 5 lít nước rồi lại đổ vào thêm 4 lít nước. Hỏi trong bể có bao nhiêu lít nước?
\shortans{7}
\loigiai{
Số lít nước có trong bể là: $8 - 5 + 4 = 7$ (lít).
\textbf{Đáp số:} $7$.
}
\end{ex}

\begin{ex}
Hiệu của số lớn nhất có một chữ số và số bé nhất có một chữ số cộng thêm 1 bằng bao nhiêu?
\shortans{10}
\loigiai{
Số lớn nhất có một chữ số là 9, số bé nhất là 0: $9 - 0 = 9$.\nCộng thêm 1: $9 + 1 = 10$.
\textbf{Đáp số:} $10$.
}
\end{ex}
